\section{Examples, Non-Examples, and Counterexamples}
\chapterline

When an exploratory proof stalls, how do you decide whether you are stuck on a technical step or whether the claim is simply false? Instead of continuing to push a proof forward, mathematicians test the statement against concrete objects to see whether it holds.

In elementary mathematics, examples are often treated as routine practice exercises after a formula has been taught. In pure mathematics, examples serve a different purpose: they are the testing ground where definitions and conjectures are investigated.

In this work, mathematical objects perform three distinct jobs:
\begin{enumerate}
\item \textbf{Exemplars}: The simplest, clearest objects satisfying a definition, confirming that the concept describes something concrete.
\item \textbf{Non-examples}: Objects that deliberately violate just one condition of a definition, showing why that specific condition is necessary.
\item \textbf{Counterexamples}: Objects that satisfy every hypothesis of a proposed statement but fail its conclusion, proving the statement false.
\end{enumerate}

\subsection{The geography of mathematical objects}

To see how these three categories interact, consider a famous historical question. A \textbf{continuous function} is one whose graph has no sudden tears or breaks, so you can sketch it without lifting your pencil from the paper. A \textbf{differentiable function} is one that is smooth enough to have a well-defined tangent line and definite slope at every point. In the early nineteenth century, mathematicians including Andr\'e-Marie Amp\`ere conjectured that every continuous function must be differentiable almost everywhere, with at most a few isolated sharp corners.

A proposed mathematical theorem of this kind has the logical form:
\[
\forall x, \quad \bigl(H(x) \implies C(x)\bigr)
\]
``For every object $x$ in our universe, if $x$ satisfies hypothesis $H$, then $x$ possesses the target property $C$.'' In the language of geometry and sets, this statement claims that the region of objects satisfying $H$ is completely contained within the region of objects possessing $C$: $H \subseteq C$.

Because every differentiable function on $\R$ is automatically continuous ($C \subseteq H$), the conjecture was really asking whether the two regions coincide: is continuity alone enough to force smoothness?

\begin{center}
\begin{tikzpicture}[
    font=\small,
    universe/.style={draw=gray!50, fill=gray!2, rounded corners=4pt, thick},
    hypo/.style={draw=sourceblue!80, fill=sourceblue!10, rounded corners=10pt, thick},
    prop/.style={draw=targetgreen!80, fill=targetgreen!15, rounded corners=8pt, thick}
]

\draw[universe] (0, 0) rectangle (12.4, 5.2);
\node[anchor=north west, font=\footnotesize\scshape, text=gray!70] at (0.2, 5.0) {Universe of Real-Valued Functions $f : \R \to \R$};

\draw[hypo] (0.6, 0.4) rectangle (8.4, 4.5);
\node[anchor=north west, font=\footnotesize\bfseries, text=sourceblue!90!black] at (0.8, 4.3) {Hypothesis $H$: Continuous Functions};

\draw[prop] (1.0, 0.7) rectangle (4.2, 3.7);
\node[anchor=north west, font=\footnotesize\bfseries, text=targetgreen!90!black] at (1.2, 3.5) {Target $C$: Differentiable};

\node[circle, fill=targetgreen!80!black, inner sep=2.5pt] (ex) at (2.6, 2.2) {};
\node[below=2pt, font=\scriptsize, align=center] at (ex) {
    \textbf{Exemplar}\\
    $f(x) = x^2$\\
    \color{gray!70!black}Smooth everywhere
};

\node[circle, fill=trapred, inner sep=2.5pt] (cx) at (6.3, 2.3) {};
\node[below=2pt, font=\scriptsize, align=center] at (cx) {
    \textbf{\color{trapred!90!black}Counterexample}\\
    $f(x) = |x|$\\
    \color{gray!70!black}Continuous, corner at $0$
};

\node[circle, fill=quotientorange!90!black, inner sep=2.5pt] (nx) at (10.4, 2.3) {};
\node[below=2pt, font=\scriptsize, align=center] at (nx) {
    \textbf{\color{quotientorange!90!black}Non-Example}\\
    $f(x) = \operatorname{sgn}(x)$\\
    \color{gray!70!black}Jump at $0$ (not in $H$)
};

\end{tikzpicture}
\end{center}

Look closely at the three distinct territories shown in the diagram:
\begin{itemize}
\item \textbf{Standard Exemplars ($f(x) = x^2$)}: Sitting securely inside the target property $C$ (and therefore inside $H$), exemplars confirm that the concept is inhabited. They provide an anchor for geometric intuition.
\item \textbf{Counterexamples ($f(x) = |x|$)}: Sitting inside the hypothesis $H$ (continuous) but outside the target $C$ (differentiable), this single function disproves the universal claim. At $x = 0$, the left-hand slope is $-1$ while the right-hand slope is $+1$, so no tangent line exists. A single function in the difference $H \setminus C$ shows the statement is false.
\item \textbf{Non-examples of the hypothesis ($f(x) = \operatorname{sgn}(x)$)}: The sign function $\operatorname{sgn}(x)$ equals $+1$ for $x > 0$, $-1$ for $x < 0$, and $0$ at the origin. Because it jumps abruptly at $0$, it is not continuous: it sits completely outside $H$. It cannot be used to test whether $H \implies C$, because it never met the entrance condition of the premise.
\end{itemize}

\begin{recoverybox}
\textbf{Pause and think}: Could a function ever sit inside $C$ (differentiable) but outside $H$ (continuous)? No! A core theorem of calculus establishes that if a function is differentiable at $x_0$, then $\lim_{x \to x_0} (f(x) - f(x_0)) = \lim_{x \to x_0} \frac{f(x)-f(x_0)}{x-x_0}(x-x_0) = f'(x_0) \cdot 0 = 0$, which forces $f$ to be continuous at $x_0$. In our geography, the region of differentiable functions outside continuous functions is strictly empty.
\end{recoverybox}

\subsection{The art of the minimal counterexample}

When testing whether a mathematical statement is true, beginners often try to construct elaborate counterexamples. In practice, however, if a statement is false, it often fails on a very simple object.

Consider the claim:
\begin{quote}
``For all real numbers $x$ and $y$, if $x^2 = y^2$, then $x = y$.''
\end{quote}
You do not need a complicated function to test this claim. Testing the simplest signed integers is enough:
\[
x = 1, \quad y = -1 \implies (1)^2 = (-1)^2 = 1, \quad \text{yet } 1 \neq -1.
\]
Testing $x = 1$ and $y = -1$ immediately disproves the claim.

Similarly, consider the statement ``Every prime number is odd.'' The minimal counterexample is the number $2$. It is the only even prime, and serves as an important counterexample in elementary number theory.

\subsubsection*{The danger of trusting early patterns}
Why is a proof required rather than twenty successful examples? Consider Leonhard Euler's famous prime-generating polynomial:
\[
P(n) = n^2 + n + 41.
\]
Let us test the first few whole numbers:
\begin{align*}
n = 0 &\implies P(0) = 41 \quad (\text{prime}), \\
n = 1 &\implies P(1) = 43 \quad (\text{prime}), \\
n = 2 &\implies P(2) = 47 \quad (\text{prime}), \\
n = 3 &\implies P(3) = 53 \quad (\text{prime}).
\end{align*}
If you continue testing by hand, you will find that for every single integer from $n = 0$ all the way to $n = 39$, the polynomial generates a prime number! Forty consecutive successes. In empirical sciences, forty consecutive confirmations with zero failures provide strong inductive evidence.

Yet at $n = 40$, the pattern fails:
\[
P(40) = 40^2 + 40 + 41 = 40(40 + 1) + 41 = 40(41) + 41 = 41(40 + 1) = 41^2 = 1681 = 41 \times 41.
\]
$1681$ is composite. Forty positive cases do not guarantee that the formula always produces primes.

The diagram below visualizes this: forty consecutive tests produce prime numbers, yet the forty-first test disproves the general claim.

\begin{center}
\begin{tikzpicture}[
    font=\small,
    streak/.style={draw=targetgreen!80, fill=targetgreen!8, rounded corners=4pt, thick, text width=12.6cm, inner sep=6pt},
    snap/.style={draw=trapred!80, fill=trapred!8, rounded corners=4pt, thick, text width=12.6cm, inner sep=6pt}
]

% Upper Box: 40 Consecutive Successes
\node[streak] (top) at (6.3, 2.2) {
    \textbf{\color{targetgreen!90!black}Forty Consecutive Primes: $n = 0, 1, 2, \ldots, 39$ \quad (\textbf{100\% Success Rate})}\\[3pt]
    \scriptsize
    $P(0) = 41$ \textbf{\color{targetgreen!80!black}\checkmark} \quad
    $P(1) = 43$ \textbf{\color{targetgreen!80!black}\checkmark} \quad
    $P(2) = 47$ \textbf{\color{targetgreen!80!black}\checkmark} \quad
    $P(3) = 53$ \textbf{\color{targetgreen!80!black}\checkmark} \quad
    $\cdots$ \quad
    $P(39) = 1601$ \textbf{\color{targetgreen!80!black}\checkmark}\\[3pt]
    \emph{Checking small values suggests a pattern, but does not prove it.}
};

% Lower Box: The Sudden Fracture
\node[snap] (bot) at (6.3, -0.6) {
    \textbf{\color{trapred!90!black}The Minimal Counterexample: $n = 40$ \quad (\textbf{Conjecture Disproved})}\\[3pt]
    \scriptsize
    $P(40) = 40^2 + 40 + 41 = 40(41) + 41 = 41(40 + 1) = 41^2 = 1681 = 41 \times 41$\\[3pt]
    \textbf{\color{trapred!90!black}Composite.} Factored by $41$. A single counterexample shows the statement is false.
};

% Connecting Label and Arrows
\node[font=\scriptsize\bfseries, text=trapred!90!black, inner sep=4pt] (mid) at (6.3, 0.8) {Pattern breaks at $n=40$};
\draw[thick, trapred!80] (top.south) -- (mid.north);
\draw[-{Latex[length=2.5mm]}, thick, trapred!80] (mid.south) -- (bot.north);

\end{tikzpicture}
\end{center}

\subsubsection*{Common test cases}
When testing a conjecture, standard test cases often expose false statements quickly:
\begin{itemize}
\item \textbf{Numbers}: $0$ (which flattens products to zero), $1$ and $-1$ (which preserve squares while flipping signs), negative numbers (which reverse inequalities when multiplied), fractions strictly between $0$ and $1$ (where powers shrink rather than grow: $(1/2)^2 = 1/4 < 1/2$), and irrational numbers like $\sqrt{2}$ (numbers that cannot be written as a fraction of integers).
\item \textbf{Sets}: the empty set $\emptyset$ (the set containing no elements, which often satisfies conditional statements vacuously), singletons $\{a\}$ (sets with exactly one member), and disjoint sets (sets that share no elements).
\item \textbf{Functions}: the constant zero function $f(x) = 0$, the identity function $f(x) = x$, absolute value $f(x) = |x|$ (which introduces a sharp corner), step functions (which have jump discontinuities), and the reciprocal $f(x) = 1/x$ (which shoots to infinity at zero).
\item \textbf{Matrices}: the zero matrix $O$, the identity matrix $I$, projection matrices (matrices satisfying $P^2 = P$, which flatten space onto a subspace), and nilpotent matrices (non-zero matrices that square to zero, such as $N = \begin{pmatrix} 0 & 1 \\ 0 & 0 \end{pmatrix}$, where $N^2 = \begin{pmatrix} 0 & 0 \\ 0 & 0 \end{pmatrix}$).
\end{itemize}

\subsection{Stress-testing theorems: removing assumptions}

A useful practice for understanding theorems is to remove one hypothesis and construct an object showing why the theorem then fails. Each hypothesis in a theorem is necessary for the conclusion to hold.

Consider the \textbf{Extreme Value Theorem}:
\begin{theorem}[Extreme Value Theorem]
Let $f : [a, b] \to \R$ be a continuous function on a closed and bounded interval. Then $f$ attains both an absolute maximum and an absolute minimum.
\end{theorem}

Before manipulating this statement, unpack the vocabulary:
\begin{itemize}
\item A \textbf{closed interval} $[a, b]$ is one that includes both of its boundary endpoints $a$ and $b$.
\item A \textbf{bounded interval} is one that has finite length, rather than stretching off to infinity.
\item A function \textbf{attains an absolute maximum} if there is an actual point $c \in [a, b]$ where $f(c) \ge f(x)$ for all $x \in [a, b]$---meaning the function takes on its maximum value at some point in the domain, rather than merely approaching an upper bound.
\end{itemize}

Notice the three distinct hypotheses: (1) closed interval, (2) bounded interval, (3) continuous function. What happens if you drop each requirement one by one?

\begin{enumerate}
\item \textbf{Drop ``closed''}: Replace $[0, 1]$ with the open interval $(0, 1)$ and choose the simple identity function $f(x) = x$. As $x$ approaches $1$ from the left, the outputs get arbitrarily close to $1$. But $1 \notin (0, 1)$. For any point $x \in (0, 1)$, the midpoint $\frac{x+1}{2}$ is also in $(0, 1)$ and gives a strictly higher value. The ceiling $1$ is the \textbf{supremum} (the least upper bound), but it is not attained by any input in $(0, 1)$. The function has no maximum value!
\item \textbf{Drop ``bounded''}: Replace $[0, 1]$ with the unbounded ray $[0, \infty)$ and keep $f(x) = x$. The function grows without bound and attains no maximum.
\item \textbf{Drop ``continuous''}: Keep the closed and bounded domain $[0, 1]$, but introduce a single jump discontinuity at the right boundary:
\[
f(x) = \begin{cases} x & \text{if } 0 \le x < 1, \\ 0 & \text{if } x = 1. \end{cases}
\]
On $[0, 1)$, the values climb steadily toward $1$. But at the very moment $x$ reaches $1$, the output plunges to $0$. The supremum $1$ is not attained.
\end{enumerate}

\begin{center}
\begin{tikzpicture}[font=\small]

% Panel 1: Dropping "closed"
\begin{scope}[xshift=0cm]
    \draw[draw=gray!30, fill=gray!2, rounded corners=4pt] (-0.8, -1.6) rectangle (4.8, 3.8);
    \node[anchor=north west, font=\footnotesize\bfseries, text=sourceblue!90!black] at (-0.6, 3.5) {Failure 1: Dropping ``closed''};
    \node[anchor=north west, font=\scriptsize, text=gray!80!black] at (-0.6, 3.0) {Domain $(0, 1)$ is open; $f(x) = x$};

    \draw[-{Latex[length=2mm]}, thick, gray!70] (-0.2, 0) -- (3.0, 0) node[right, font=\scriptsize] {$x$};
    \draw[-{Latex[length=2mm]}, thick, gray!70] (0, -0.2) -- (0, 2.2);
    \node[left=1pt, font=\scriptsize, text=gray!70!black] at (0, 2.1) {$y$};

    \draw[dashed, trapred!80, thick] (0, 1.8) -- (2.0, 1.8);
    \node[left, font=\scriptsize, text=trapred!90!black] at (0, 1.8) {$1$};
    \node[anchor=south west, font=\tiny\bfseries, text=trapred!90!black] at (0.05, 1.85) {Supremum $= 1$};

    \draw[very thick, sourceblue] (0.06, 0.06) -- (1.74, 1.74);
    \draw[thick, sourceblue, fill=white] (0, 0) circle (2pt);
    \draw[thick, trapred, fill=white] (1.8, 1.8) circle (2.5pt);

    \draw[thick, gray!70] (1.8, 1.5pt) -- (1.8, -1.5pt) node[below=2pt, font=\scriptsize] {$1$};

    \node[below, font=\scriptsize, text=gray!90!black, align=center] at (1.8, -0.7) {
        Endpoint $1$ excluded $\implies$ peak never reached\\[2pt]
        \textbf{\color{trapred!90!black}No maximum attained!}
    };
\end{scope}

% Panel 2: Dropping "continuous"
\begin{scope}[xshift=6.0cm]
    \draw[draw=gray!30, fill=gray!2, rounded corners=4pt] (-0.8, -1.6) rectangle (4.8, 3.8);
    \node[anchor=north west, font=\footnotesize\bfseries, text=targetgreen!90!black] at (-0.6, 3.5) {Failure 2: Dropping ``continuous''};
    \node[anchor=north west, font=\scriptsize, text=gray!80!black] at (-0.6, 3.0) {Domain $[0, 1]$ is closed; jump at $x = 1$};

    \draw[-{Latex[length=2mm]}, thick, gray!70] (-0.2, 0) -- (3.0, 0) node[right, font=\scriptsize] {$x$};
    \draw[-{Latex[length=2mm]}, thick, gray!70] (0, -0.2) -- (0, 2.2);
    \node[left=1pt, font=\scriptsize, text=gray!70!black] at (0, 2.1) {$y$};

    \draw[dashed, trapred!80, thick] (0, 1.8) -- (2.0, 1.8);
    \node[left, font=\scriptsize, text=trapred!90!black] at (0, 1.8) {$1$};

    \draw[very thick, targetgreen] (0, 0) -- (1.74, 1.74);
    \fill[targetgreen] (0, 0) circle (2pt);
    \draw[thick, trapred, fill=white] (1.8, 1.8) circle (2.5pt);

    \draw[dotted, thick, trapred!70] (1.8, 1.7) -- (1.8, 0.1);
    \fill[trapred] (1.8, 0) circle (2.5pt);
    \node[right=2pt, font=\tiny\bfseries, text=trapred!90!black] at (1.8, 0.4) {$f(1) = 0$};

    \node[below=3pt, font=\scriptsize, text=gray!70] at (1.8, 0) {$1$};

    \node[below, font=\scriptsize, text=gray!90!black, align=center] at (1.8, -0.7) {
        Values climb to $1$ but drop at end\\[2pt]
        \textbf{\color{trapred!90!black}No maximum attained!}
    };
\end{scope}

\end{tikzpicture}
\end{center}

Each hypothesis was necessary: removing any one of them allows the conclusion to fail.

\begin{historybox}
In 1752, Leonhard Euler announced a geometric law: for any \textbf{polyhedron}---a solid three-dimensional shape with flat polygonal faces, straight edges, and sharp corner vertices, such as a cube or pyramid---the number of vertices ($V$), edges ($E$), and faces ($F$) satisfies:
\[
V - E + F = 2.
\]
For a cube, $8 - 12 + 6 = 2$; for a tetrahedron, $4 - 6 + 4 = 2$. In 1813, Augustin-Louis Cauchy published what appeared to be an airtight proof by removing one face and flattening the remaining network of edges onto a flat two-dimensional plane.

Yet over the following century, unexpected geometric ``monsters'' appeared. Johannes Kepler had already described self-intersecting star polyhedra (solids whose faces pass through one another), which yield $12 - 30 + 12 = -6$. In 1813, the Swiss mathematician Louis Lhuilier demonstrated that a picture-frame polyhedron with a central square tunnel gives $V - E + F = 0$, while a hollow cube floating inside another cube gives $V - E + F = 4$.

In his 1976 masterwork \emph{Proofs and Refutations}, the philosopher Imre Lakatos analyzed this history to show that mathematical discovery does not proceed in a straight line from unquestioned axioms to permanent theorems. Instead, it advances through a living dialogue of bold conjectures and surprising counterexamples. Early mathematicians reacted with ``monster-barring''---declaring that solids with tunnels or self-intersections were ``not real polyhedra'' to defend Euler's original equation. The true breakthrough came from inspecting where Cauchy's proof had quietly made an unstated assumption: Cauchy had assumed that every polyhedron could be inflated like a balloon into a round sphere.

By incorporating the missing topological condition---the \textbf{genus} $g$, which counts the number of donut-like tunnels passing through the solid---mathematicians established the universal formula:
\[
V - E + F = 2 - 2g.
\]
This quantity is the simplest example of a \textbf{topological invariant}: a numerical quantity that does not change when a shape is continuously stretched, bent, or compressed without tearing. Counterexamples led directly to the formulation of modern topological invariants.
\end{historybox}

The diagram below shows how a geometric counterexample reshapes a naive theorem: Lhuilier's tunnel polyhedron breaks Euler's original $V - E + F = 2$, revealing the hidden topological invariant $g$.

\begin{center}
\begin{tikzpicture}[
    font=\small,
    panel/.style={draw=gray!40, fill=gray!2, rounded corners=5pt, thick, inner sep=6pt}
]

% Panel 1: Convex Polyhedron (Cube, g=0)
\begin{scope}[xshift=0cm]
    \node[panel, minimum width=6.2cm, minimum height=5.2cm] (p1) at (2.6, 1.8) {};
    \node[anchor=north, font=\footnotesize\bfseries, text=sourceblue!90!black] at (2.6, 4.2) {Standard Cube: Genus $g = 0$};

    % Cube 3D drawing
    % Back hidden lines
    \draw[dashed, gray!60, thick] (1.7, 1.5) -- (3.5, 1.5);
    \draw[dashed, gray!60, thick] (1.7, 1.5) -- (1.7, 3.1);
    \draw[dashed, gray!60, thick] (1.7, 1.5) -- (0.9, 0.8);

    % Faces shading
    \fill[sourceblue!15, opacity=0.8] (0.9, 0.8) -- (2.7, 0.8) -- (2.7, 2.4) -- (0.9, 2.4) -- cycle;
    \fill[sourceblue!25, opacity=0.8] (2.7, 0.8) -- (3.5, 1.5) -- (3.5, 3.1) -- (2.7, 2.4) -- cycle;
    \fill[sourceblue!35, opacity=0.8] (0.9, 2.4) -- (2.7, 2.4) -- (3.5, 3.1) -- (1.7, 3.1) -- cycle;

    % Visible edges
    \draw[thick, sourceblue!80!black] (0.9, 0.8) -- (2.7, 0.8) -- (2.7, 2.4) -- (0.9, 2.4) -- cycle;
    \draw[thick, sourceblue!80!black] (2.7, 0.8) -- (3.5, 1.5) -- (3.5, 3.1) -- (2.7, 2.4);
    \draw[thick, sourceblue!80!black] (0.9, 2.4) -- (1.7, 3.1) -- (3.5, 3.1);

    % Vertices dots
    \foreach \coord in {(0.9,0.8), (2.7,0.8), (2.7,2.4), (0.9,2.4), (3.5,1.5), (3.5,3.1), (1.7,3.1)} {
        \fill[sourceblue!90!black] \coord circle (2pt);
    }
    \fill[gray!50] (1.7, 1.5) circle (1.8pt);

    \node[anchor=south, font=\scriptsize, align=center] at (2.6, -0.6) {
        $V = 8, \quad E = 12, \quad F = 6$\\[3pt]
        $V - E + F = 8 - 12 + 6 = \mathbf{2}$\\[2pt]
        \textbf{\color{targetgreen!80!black}Matches Euler ($2 - 2(0) = 2$)}
    };
\end{scope}

% Panel 2: Tunnel Polyhedron (Picture frame, g=1)
\begin{scope}[xshift=6.8cm]
    \node[panel, minimum width=6.2cm, minimum height=5.2cm] (p2) at (2.6, 1.8) {};
    \node[anchor=north, font=\footnotesize\bfseries, text=trapred!90!black] at (2.6, 4.2) {Tunnel Solid: Genus $g = 1$};

    % Schematic frame representation
    % Outer block
    \draw[thick, trapred!80!black, fill=trapred!12] (0.7, 0.8) rectangle (4.5, 3.2);
    % Inner tunnel
    \draw[thick, trapred!80!black, fill=white] (1.7, 1.4) rectangle (3.5, 2.6);
    \node[font=\tiny\bfseries, text=trapred!90!black] at (2.6, 2.0) {Square Tunnel};

    % 3D frame depth lines
    \draw[thick, trapred!70!black] (4.5, 0.8) -- (4.9, 1.3) -- (4.9, 3.7) -- (4.5, 3.2);
    \draw[thick, trapred!70!black] (4.9, 3.7) -- (1.1, 3.7) -- (0.7, 3.2);
    \draw[thick, trapred!70!black] (3.5, 2.6) -- (3.8, 2.9);

    \node[anchor=south, font=\scriptsize, align=center] at (2.6, -0.6) {
        $V = 16, \quad E = 32, \quad F = 16$\\[3pt]
        $V - E + F = 16 - 32 + 16 = \mathbf{0}$\\[2pt]
        \textbf{\color{trapred!90!black}Breaks naive formula! ($2 - 2(1) = 0$)}
    };
\end{scope}

\end{tikzpicture}
\end{center}

\subsection{Why set theory provides the shared framework}

Seeking exemplars, non-examples, and counterexamples grounds abstract statements in concrete behavior. By testing extreme cases ($0, 1, -1, \emptyset$) and removing hypotheses one by one, you discover why each condition in a theorem is needed.

Testing statements across numbers, functions, and shapes raises a natural question: what common language connects these different kinds of mathematical objects? In Chapter~10 (\emph{Sets: The Grammar of Modern Mathematics}), we introduce set theory, the framework used throughout modern mathematics to define and organize objects.

\subsection{Exercises}

\paragraph{Quick mental checks.} For each statement below, find a counterexample:
\begin{enumerate}
\item For all real numbers $a$ and $b$, if $ab = 0$, then $a = 0$ and $b = 0$.
\item For every real number $x$, if $x^2 > 4$, then $x > 2$.
\item If a polygon in the plane has four equal sides, it must be a square.
\item For all positive integers $a, b, c$, if $a$ divides the product $bc$, then $a$ divides $b$ or $a$ divides $c$.
\end{enumerate}

\paragraph{Routine practice.} Provide a concrete counterexample for each false claim and identify precisely which condition or hypothesis fails:
\begin{enumerate}[resume]
\item For all sets $A, B, C$, if $A \cup B = A \cup C$, then $B = C$.
\item For all sets $A, B, C$, if $A \cap B = A \cap C$, then $B = C$.
\item If a differentiable function $f : D \to \R$ satisfies $f'(x) = 0$ for every $x \in D$, then $f$ is a constant function. (Hint: What if the domain $D$ consists of two disconnected pieces, such as $(-\infty, 0) \cup (0, \infty)$?)
\item Every sequence of rational numbers that converges in $\R$ must have a rational limit. (Hint: Consider decimal truncations of an irrational number, such as $1, 1.4, 1.41, 1.414, \ldots \to \sqrt{2}$.)
\item For all real numbers $x$, if $x^2 \ge x$, then $x \ge 1$.
\end{enumerate}

\paragraph{Deeper exploration.}
\begin{enumerate}[resume]
\item \textbf{Euler's sum-of-powers conjecture}: In 1769, Leonhard Euler conjectured that for all integers $n > 2$, it requires at least $n$ positive $n$-th powers to sum to another $n$-th power. For $n = 3$, this asserts that two cubes cannot sum to a cube (the $n=3$ case of Fermat's Last Theorem), requiring three cubes ($3^3 + 4^3 + 5^3 = 6^3$). For $n = 5$, Euler's conjecture claimed that no four fifth powers could sum to a fifth power: $a^5 + b^5 + c^5 + d^5 = e^5$ was conjectured to be impossible. In 1966, L.~J.~Lander and T.~R.~Parkin discovered the following equality:
\[
27^5 + 84^5 + 110^5 + 133^5 = 144^5.
\]
Verify by direct calculation that this identity holds. Explain why this single numerical counterexample permanently settled a conjecture that had stood for 197 years.
\item \textbf{Patterns and proof}: In 1650, Pierre de Fermat inspected the numbers $F_n = 2^{2^n} + 1$:
\[
F_0 = 3, \quad F_1 = 5, \quad F_2 = 17, \quad F_3 = 257, \quad F_4 = 65537.
\]
Observing that all five numbers are prime, Fermat conjectured that $F_n$ is prime for every non-negative integer $n$. In 1732, Leonhard Euler factored the very next term:
\[
F_5 = 2^{32} + 1 = 4{,}294{,}967{,}297 = 641 \times 6{,}700{,}417.
\]
Explain the philosophical difference between scientific induction (generalizing from observed instances) and mathematical proof: why can no finite collection of examples ever prove a statement of the form $\forall n \in \N, \, P(n)$?
\end{enumerate}

\newpage
